\section{確率分布 : 形を言葉にし, 推論の足場をつくる}

  \subsection{問い : ヒストグラムの"形"は, 偶然か, 仕組みか}
    前章で図にしてみると, いくつかの「形」が目に入ったはずだ.
    山がどこにあるか, どちら側に裾が長いか, ぽつんと離れた点はあるか.

    \begin{itemize}
      \item \textbf{この形はたまたまか}. もし同じ条件で\textit{もう一度}データを集めたら, 山の位置や裾の感じは同じになるのか.
      \item \textbf{何が形を作っているのか}. ジャンルの違い, 公開時期, おすすめ表示など\textit{仕組み}が形に現れているのか, それとも小さな偶然の積み重ねなのか.
      \item \textbf{どこまでを"言える"とみなすか}. 今の図の印象を, 全体の話としてどこまで持ち出してよいのか.
    \end{itemize}

    たとえば, 20作品の視聴回数のヒストグラムで右に長い裾が見えたとしよう.
    それは「人気作品への集中」という\textit{仕組み}が作った形かもしれない.
    あるいは, たまたま今月バズった作品が混ざって\textit{偶然}そう見えただけかもしれない.
    ビン幅を少し動かすだけで山の数が変わることもあるし, サンプルが少ないと外れ値1つで印象が大きく変わる.

    では, どう判断するか. まずは, 図を見て形を言葉にする素朴なやり方から始めてみよう.

  \subsection{素朴な方法と限界 : 「形を見る」だけでは足りない理由}
    前章のやり方, すなわちヒストグラムを描いて「右裾が長い」「山が二つある」と目で読み取る方法は, 出発点として有効だ.
    しかし, 困ったことが二つある.

    \begin{enumerate}
      \item \textbf{形の不安定さ}. ビン幅を変えるだけで山の数が変わる. サンプルサイズが小さいと, 同じ条件で取り直したとき形がまるで別物になることもある.
            つまり, 今の形のどこまでが「仕組み」で, どこからが「偶然」なのかを, 目で見るだけでは切り分けられない.
      \item \textbf{要約値の一点語り}. 「平均視聴回数は 4{,}500 回」と言い切っても, 取り直したら 3{,}800 回かもしれない.
            一点の数字に\textbf{どのくらいの幅がつくのか}が分からなければ, その数字をどこまで信じてよいか判断できない.
    \end{enumerate}

    どちらの困りごとも, 根っこは同じだ.
    \textbf{「取り直したらどれくらいばらつくか」が分からない}のである.
    このばらつきを見積もるための道具が, 次に見る\textbf{確率分布}と, そこから導かれる\textbf{標準誤差}（SE）である.

  \subsection{なぜ理論分布を持ち出し, 統計量のばらつきを意識するのか}
    前章で見えたのは\textbf{形}だった.
    山の位置や裾の長さは, たしかに多くを語る.
    しかし, それだけでは足りない.
    所見を\textbf{言葉として共有}し, なおかつ\textbf{過大に語らずに一般化}するには, もう一歩だけ道具立てが要る.
    その道具立てが\textbf{理論分布}と, \textbf{統計量を同じ条件で繰り返し計算したときの分布}である.

    理論分布に\textbf{当てはめること自体が目的ではない}.
    \begin{enumerate}
      \item \textbf{まず, 手元の形を見る}. ヒストグラムや箱ひげ図から, 対称性, 裾, 多峰性, 外れの有無を素直に読む. これが\textbf{作法}（平均で語るか中央値か, 変換や外れ値処理が要るか）を決める.
      \item \textbf{次に, 形を言葉にする}. 理論分布は\textbf{モデルの形}を与える\textbf{言語}だ.
            「釣鐘型に近い」「右に裾が長い」「回数データでばらつきがこうだ」などと短く言えるだけで, 所見は他者と\textbf{比較可能}になる.
            データを型に押し込むためではない.
      \item \textbf{そして, ばらつきを意識する}. 同じ条件で取り直したら, 平均や割合などの\textbf{統計量}はどのくらい\textbf{ばらつく}か.
            この「取り直したときの並び」を思い描くのが\textbf{統計量を同じ条件で繰り返し計算したときの分布}である.
            後章で扱う\textbf{幅づけ}や\textbf{珍しさの判断}は, この分布を\textbf{物差し}にして行う.
    \end{enumerate}

    理論分布がないと, どう困るか.
    \begin{itemize}
      \item \textbf{形の語彙と基準を与える}. 正規, 二項, ポアソン, 指数などは, 形を\textbf{名前で呼べる}ようにし, 観察結果を\textbf{共有可能な基準}に載せる. ぴったり一致を求めないからこそ, 語彙として使える.
      \item \textbf{仕組みの仮説を短く描ける}. たとえば「小さな寄与の和が結果を作る」なら釣鐘型に近づきやすい, 「一定率で事象が起きる回数」ならポアソンが手がかりになる.
            これは\textbf{現象のモデル}を素描するためであり, 現実をそのまま置き換えるためではない.
      \item \textbf{使える道具と注意点が見える}. 形が分かれば, どの要約や検定が\textbf{前提に合うか}が分かり, 逆に合わないときは\textbf{変換や頑健な指標}に切り替える判断ができる.
    \end{itemize}

    私たちの推定値は一度きりの一本の数ではない.
    「取り直したら取り得た多くの値」の中の一点にすぎないのだ.
    ばらつきの\textbf{大きさと形}が分かってはじめて, どこまで一般化してよいかを語れる.

    \medskip
    \noindent 最小限の用語を一言でそろえる.
    \textbf{理論分布}＝現象を説明・比較するための\textbf{モデルの形}.
    \textbf{統計量を同じ条件で繰り返し計算したときの分布}＝同条件で取り直したときの\textbf{統計量のばらつき方}.
    出発点はつねに\textbf{手元の形}である.
    この二つを区別することで, 手元の形を語る言葉と, 推定値の不確かさを測る物差しが手に入る.

    この視点が足場になる.
    次に見る\textbf{正規分布と中心極限定理}で, とりわけ\textbf{平均のばらつき方が整ってくる理由}が見えてくる.

    \subsubsection*{正規分布 : 多数の小さな変動の和}
      1つの測定値は, 多くの小さな要因の足し合わせで決まることが多い（製造の微差, 測定の誤差, 日々の変動など）.
      こうした「小さな変動の和」は, 釣鐘型の形を生みやすい.

    \subsubsection*{中心極限定理（芯だけ）}
      同じ集団から大きさ $n$ の標本を何度も取り直し, \textbf{平均}を毎回計算する.
      $n$ を大きくしていくと, その平均の分布は\textbf{正規分布}に近づく.
      元の形は問わない.
      前提はただ一つ, 各観測が\textbf{互いに独立で同じ条件から得られていること}（独立同分布, i.i.d.）だ.
      この条件さえ満たせば, 元の形が多少ゆがんでいても, 平均は\textbf{釣鐘型}に落ち着いてくる.
      ただし, ゆがみが強いときは $n$ が大きくないと近づきが遅い. これは前提ではなく\textbf{収束の速さ}の問題である.
      だから本書で使う「推定値 ± およそ $2\times$ \textbf{標準誤差（SE）}」がよく効く.

    \subsubsection*{大数の法則との違い（ひと言）}
      大数の法則は「平均が真の値に近づく」話.
      中心極限定理は「\textbf{近づき方の形}が釣鐘型になる」話. どちらも, 幅で語る根拠になる.

    \subsubsection*{データの形と統計量のばらつき}
      実データの形は\textbf{理論分布に完全には合わない}ことが多い. 視聴回数は右に長い, 待ち時間は裾が太い, など.
      ゆえに, 図で確かめ, 必要なら\textbf{変換}（例: 対数）や\textbf{頑健な指標}（中央値, IQR）を使う.

      では, データの形が合わないのに, なぜ統計量のばらつきは理論で扱えるのか.
      平均, 割合, 差, 比といった\textbf{統計量のばらつき方}は, 条件が整えば\textbf{理論分布}に落ち着くからだ.
      平均のばらつきは\textbf{正規}に近づき, 割合や差のばらつきも\textbf{式で広がりが計算}できる.
      \textit{要点}：\underline{データを無理に理論分布へ合わせる必要はない}. 必要なのは, \underline{統計量のばらつき方}をつかんで, \underline{幅}を判断すること.

    \medskip
    後の章で, 平均の幅づけには $t$ 分布を, ばらつきの大きさには $\chi^2$ 分布を, 二つのばらつきの比には $F$ 分布を物差しとして使う.
    これらは\textbf{データそのものの型}ではなく, \textbf{統計量がどれくらいなら珍しいか}を測る道具である. 詳しくは推定・検定の章で手を動かしながら確かめる.

    \begin{mybox}{代表的な理論分布 : 仕組みから形が決まる}
    \textbf{一様分布（Uniform）}.
    区間 $[a,b]$ のどの値も同じ確からしさで現れる状況. 初期化の乱数, 無作為抽選の理想像.
    形は平坦. 平均は $(a+b)/2$, 広がりは $(b-a)^2/12$.
    目印はヒストグラムが山なしで平ら, 端で急に切れること.
    現実にぴったり平坦はまれ（丸め, 測定下限・上限の影響が出やすい）.

    \medskip
    \textbf{二項分布（Binomial）}.
    「成功/失敗」を確率 $p$ で独立に $n$ 回繰り返し, 成功回数を数える. 例：クリックするか, 高評価をつけるか.
    $np(1-p)$ が大きいと釣鐘に近づき, $p$ が端に近いと片側に歪む. 平均 $np$, 広がり $np(1-p)$.
    目印は $n$ 回中の成功回数として集計されたデータ.
    試行が独立でない, $p$ がばらつくと合わない（混合や過分散に注意）.

    \medskip
    \textbf{ポアソン分布（Poisson）}.
    一定の平均率 $\lambda$ で, 独立にまばらに起こる出来事の区間内の回数. 例：1日あたりのレビュー投稿件数.
    右裾が長い. 平均と分散がともに $\lambda$.
    目印は平均と分散がほぼ等しいこと.
    分散が平均より大きいなら過分散（負の二項などを検討）.

    \medskip
    \textbf{指数分布（Exponential）}.
    ポアソン過程で次の出来事までの待ち時間.
    右に長い尾, すぐ減衰. 平均 $1/\lambda$. 記憶なしの性質（過去の長さが次に影響しない）.
    目印は生存曲線が対数軸で直線.
    発生率が時間で変わると合わない（ワイブルなどへ）.

    \medskip
    \textbf{正規分布（Normal）}.
    多数の小さな寄与の和. 測定誤差や寸法のばらつき.
    対称の釣鐘. 平均と分散で形が決まる.
    外れ値や重い尾に弱い（中央値や IQR 併用, 変換を検討）.

    \medskip
    \textbf{対数正規分布（Lognormal）}.
    掛け算の積み重ねで大きさが決まる. 例：視聴回数, 所得の一部.
    正の値のみ, 右に長い尾. $\log$ を取ると正規に近づく.
    目印は $\log$ 変換後のヒストグラムが釣鐘.
    平均が大きく外れ値に引っ張られる. 中央値の併記が有用.

    \medskip
    \noindent \textit{選び方の目安}：\textbf{仕組み（数え上げか, 待ち時間か, 足し算か掛け算か）}と\textbf{形（対称か, 裾の長さ）}をセットで見る.
    ぴったり一致を目指すのではなく, \textbf{言語と基準}として使う.
    \end{mybox}

    道具がそろったところで, 実際のデータに当てはめる手順に入る.

  \subsection{形を確かめ, ばらつきに幅をつける}
    \noindent\textbf{前提}：データの各観測が互いに独立で, 同じ条件から得られていること（i.i.d.）. 極端な外れ値や強い歪みがあるときは, 変換や頑健な指標の併用を検討する. 似た観測が連続する, 期間が混ざるなどで独立性が崩れるときは, 分析者が切り取り基準を明示する.

    \begin{itemize}
      \item \textbf{形を確認}. ヒストグラムと\textbf{密度曲線}を重ねる.
            山の数, 歪み, 裾の長さを一文で要約する.
      \item \textbf{仕組みを当てる}. \textbf{数え上げ}なら二項/ポアソン, \textbf{待ち時間}なら指数, \textbf{掛け算で増える}なら対数正規を\textbf{候補語}として使う.
            ぴったり一致は狙わず, \textbf{説明の言語}として用いる.
      \item \textbf{ばらつきを意識}. 平均や割合など, 使う統計量ごとに\textbf{SE}の見積もりを添え, 「推定値 ± 約 $2SE$」で\textbf{幅}を示す（根拠は CLT と正規近似）.
    \end{itemize}

    \noindent\textit{注意}：実データが完全に正規, はまれだ. 平均・分散は外れ値に弱く, 歪みが強いときは中央値・分位や変換を併用する. 「ぴったり」は求めず, 図で形を見て道具を選ぶ.

  \subsection{実践 : 20作品データで体験する}
    \begin{enumerate}
      \item \textbf{視聴回数の形}. ヒストグラムと密度曲線. 右裾が長ければ, 対数を取って再描画. \textbf{釣鐘に近づく}か観察する.
      \item \textbf{レビュー数の当てどころ}. 平均と分散を比較し, 分散が平均に近いかを確認（ポアソン的か\textbf{目安}を取る）.
      \item \textbf{平均を繰り返し計算したときの分布の感触}. 視聴回数の\textbf{平均}の SE を計算し, 「推定値 ± 約 $2SE$」の\textbf{幅}を併記. 図の印象と照合する.
      \item \textbf{割合でも同様に}. 星評価 $\ge 4.0$ の\textbf{割合}について, 形の所見と幅の提示を行う（端に近いと非対称に注意）.
    \end{enumerate}
    \textit{所見例}：「視聴回数は右裾が長い. 対数を取るとほぼ釣鐘型になる. 平均の SE は $\ldots$ で, 幅（± 約 $2SE$）を併記する.」

    \textit{記録の約束}：軸・単位・期間, 切り取り基準, 欠損の扱いを明示する.

  \subsection{結果と次の一手 : 幅をもって主張する道が開けた}
    分布を知るとは, 図の形を\textbf{言葉にする}ことだ.
    実データの形は完璧には合わないが, \textbf{統計量のばらつき方}は理論形に\textbf{落ち着く}.
    ばらつきの大きさが SE として手に入ったことで, 「推定値 ± 約 $2SE$」と\textbf{幅をもって主張する}道が開けた.
    次章では, この幅を使い, \textbf{点推定と信頼区間}として数値の不確かさを正面から扱う.
